% s02-sec1.tex —— §1 L 函数的极点与函子性（原文 p.3–12）
\section{\texorpdfstring{$L$ 函数的极点与函子性}{L 函数的极点与函子性}}
\srcpage{3}

关于自守表示所附 $L$ 函数的基本定义，我们请读者参阅《自守形式、表示与 $L$ 函数》[BC]。关于自守表示理论还有若干更新的优秀导论，例如 [AEK]，但它们不处理 $L$ 函数。读者在阅读本文时需要大致理解稳定迹公式及其预期应用，因而需要理解内窥之外的自守形式。在某些方面，尽管 Arthur 做出了卓绝而根本的贡献，稳定迹公式仍然只以粗略的形式存在 [A2, A3, K2, L2]。我们希望在本文及后续文章中说明它的目标可能是什么，以及如何借助它有望实现这些目标。我们只是摸索着前进，但为了给读者指路，我们先描述背景。

$L$ 函数是（对所有有限或无限位 $v$ 取）乘积
\begin{equation}\tag{1.1}
L(s,\pi,\rho)=\prod_v L(s,\pi_v,\rho),
\end{equation}
这里的 $L$ 函数附属于一个群 $G$ 的自守形式——更确切地说，附属于自守形式的 $L$-包——以及其 $L$-群的一个表示 $\rho$。

Arthur（[A1]）提出了自守形式的一种分类，这种分类只有借助函子性才能实现，而函子性本身无疑也只能与分类同时建立。要理解我们的策略，至少必须理解这种分类的要素会有哪些。为了一般地建立它，人们预期要使用迹公式与递推。在接下来几页中我们将看到，如何在操作这一公式时顾及分类的推论。

人们预期可以把若干对象联系到一个自守表示 $\pi=\pi_G$：首先是一个同态 $\varphi=\varphi_\pi:\operatorname{SL}(2)\to{}^LG$；其次，对几乎所有的位 $v$，是 $\varphi(\operatorname{SL}(2))$ 在 ${}^LG_v$ 中的连通中心化子 ${}^{\lambda}G_{\varphi,v}$ 里的一个共轭类 $\{A^G(\pi_v)\}=\{A(\pi_v)\}$。\yizhu{原文把 $\varphi(\operatorname{SL}(2))$ 在 ${}^LG_v$ 中的连通中心化子记作 ${}^{\lambda}G_{\varphi,v}$，把假设的同态 $\xi$ 的像的 Zariski 闭包记作 ${}^{\lambda}H_\pi$（见下页）。「$\lambda$」是原文记号的一部分，不是指数，也不表示某个群。}{依据：原书 p.3–4}

\srcpage{4}
甚至对每个位 $v$，还有局部 Weil 群到 ${}^{\lambda}G_{\varphi,v}$ 中的一个同态。人们希望构造一个\textbf{自守 Galois 群}，以及从这个群到整体中心化子 ${}^{\lambda}G_\varphi$ 的一个同态 $\xi$，使得它生成这些局部类 $\{A(\pi_v)\}$。这是一个计划，而且是一个理应导致自身实现的计划。\yizhu{「自守 Galois 群」是作者对未来理论的设想：一个其局部同态恰好生成各 $\{A(\pi_v)\}$ 的 Galois 型群。原文随即指出，这样的群会「大得无用」；真正有用的是 ${}^{\lambda}H_\pi$。}{依据：原书 p.4}

事实上，这样的 Galois 群会大得无用且笨重。更有用的是由假设的映射 $\xi$ 的像的 Zariski 闭包所定义的群 ${}^{\lambda}H_\pi$。它们比同态 $\xi$ 更为根本；人们希望，在确定了这些 ${}^{\lambda}H_\pi$ 之后，能够定义一个自守 Galois 群。确定 ${}^{\lambda}H_\pi$ 的问题在 [L3] 中已经开始处理。人们甚至要在 $F$ 上找一个约化群 $H$、一个核为中心的满同态
\[
\psi:{}^LH\longrightarrow{}^{\lambda}H\subset{}^{\lambda}G_\varphi,
\]
以及一个表示 $\pi_H$，使得对几乎所有的 $v$ 有
\[
\{A(\pi_v)\}=\{\psi(A(\pi_{H,v}))\}.
\]
按照 Arthur 提出的分类，有一类特殊的表示，称为\textbf{Ramanujan 型}，它们满足 Ramanujan 猜想的结论。\yizhu{Ramanujan 型是 Arthur 分类中的术语：其局部参数在每个位都是酉的，从而 $L$ 函数在临界线右侧没有极点。原文接着把这一性质表述为「类 $\{A(\pi_{H,v})\}$ 是酉的」。}{依据：原书 p.4；并见 [A1]} 我们希望 $\pi_H$ 是 Ramanujan 型，也就是说，这些类 $\{A(\pi_{H,v})\}$ 是酉的。在 [L3] 中处理的问题是：在迹公式的框架内，由 $\pi$ 构造出数据 $\varphi$、$H$、$\pi_H$ 与 $\psi$。我们顺便指出，并不预期这些数据是唯一的。会有与这些构造相关的重数问题，Song Wang [WS] 已具体处理过。

人们预期能够利用稳定迹公式做适当的递推，并依据函子性原则证明两件事：一方面，类 $\{A(\pi_{H,v})\}$ 是酉的，从而对 ${}^LH$ 的每个表示 $\rho_H$，$\rho_H(A(\pi_{H,v}))$ 的特征值绝对值都为 $1$；另一方面，$L$ 函数 $L(s,\pi_H,\rho_H)$ 在 $\operatorname{Re}s>1$ 上全纯，在直线 $\operatorname{Re}(s)=1$ 上至多有有限个极点，且在这条临界线上没有零点。（若 $F$ 是 $q$ 元有限域上的函数域，则按模 $2\pi i/\ln q$ 计算极点。）我们注意到，Euler 乘积中有限个位的影响极小，因此所描述的行为归根结底是部分 $L$ 函数的行为。

\srcpage{5}
设 $S$ 为有限位集，定义部分 $L$ 函数
\begin{equation}\tag{1.1'}
L_S(s,\pi,\rho)=\prod_{v\notin S}L(s,\pi_v,\rho).
\end{equation}
原文此处把 (1.1) 与部分 $L$ 函数共用同一编号，译本为免混淆，将部分 $L$ 函数标作 (1.1$'$)，其余编号照原文。
表示 $\rho$ 与同态 $\varphi\times\psi$ 定义了乘积 $\operatorname{SL}(2)\times{}^LH$ 的一个表示，它分解为直和
\begin{equation}\tag{1.2}
\bigoplus_j\sigma_j\otimes\rho^j_H,
\end{equation}
其中每个 $\sigma_j$ 不可约。设 $\operatorname{SL}(2)$ 的表示 $\sigma_j$ 的维数为 $m_j+1$，因而其最高权为 $\mu_j:(1,-1)\mapsto m_j$。则 $L$ 函数分解为乘积
\begin{equation}\tag{1.3}
L_S(s,\pi,\rho)=\prod_j\prod_i L_S(s+i,\pi_H,\rho^j_H),
\qquad i\in\left\{\frac{m_j}{2},\frac{m_j}{2}-1,\dots,-\frac{m_j}{2}\right\}.
\end{equation}

显然，这些函数的行为强烈依赖于整数 $m_j$。当 $m_j=0$ 时，临界线是 $s=1$，临界带是 $0\le\operatorname{Re}s\le1$，其中心为 $\operatorname{Re}s=1/2$。当 $m_j>0$ 时，奇点随着 $s$ 减小而更早出现。然而若 $m_j>0$，则 $\dim H<\dim G$，于是我们可以对 $\dim G$ 或 $\dim\rho$ 作归纳，假定我们理解 $L_S(s,\pi_H,\rho^j_H)$ 的行为，并且已经对 $H$ 证明了 Arthur 的假设，特别是 Ramanujan 假设。我们提醒，后一假设等价于：对于 $\varphi$ 平凡的表示，$L(s,\pi,\rho)$ 可解析延拓到直线 $\operatorname{Re}s=1$。

更确切地说，从迹公式首先得到和
\begin{equation}\tag{1.4}
\sum_\pi\prod_{v\in S}\operatorname{tr}(\pi_v(f_v))L_S(s,\pi,\rho)
\end{equation}
的公式。函数 $f_v$ 光滑且紧支，大致任意，使我们能在论证末尾把表示 $\pi$——至少是它们的稳定类——彼此分离。此外，这类表达式隐含地对 $S$ 外非分歧的 $\pi$ 求和。既然我们已知子群 ${}^LH$ 以及使 $\varphi$ 的像与 ${}^LH$ 交换的 $\varphi$，我们就能通过归纳，从群 $H$ 的迹公式推出 (1.4) 中来自不等于 $G$ 的 $H$ 的贡献的公式。

\srcpage{6}
然而必须考虑到一个 $\pi$ 可能来自多个 $H$，例如来自满足 ${}^LH_1\subset{}^LH_2$ 的 $H_1$ 与 $H_2$。例子还提示，$\pi$ 可能来自两个同构但不共轭的子群，这一现象似乎与大于 $1$ 的重数的存在有关 [WS]。随着研究进展，人们期待会有澄清。于是，出现在差中的每个 $\pi$ 都假定是 Ramanujan 型，因此这个差应期望是 $\operatorname{Re}s>1$ 中的全纯函数；主要问题将是证明这一点并由此推出 Ramanujan 猜想。显然这个问题不会容易。我们在 §5 中从一个容易的侧面入手。结果并非没有意义！显然，上述讨论隐含了一般函子性的证明。

最好分两步排除不同于 $G$ 的 $H$。首先减去满足 $\varphi$ 非平凡的配对 $(\varphi,\psi)$ 的贡献。对它们，(1.3) 中函数的极点出现在更大的 $\operatorname{Re}s$ 处，因此可以在紧靠极点右侧的小区间内依次、独立地进行比较。目标是从右到左依次处理。最后我们就把 $\varphi$ 假定为平凡的 $\pi$ 分离出来，并证明 Arthur 假设所预言的函子性至少对非平凡 $\varphi$ 成立。使 $\varphi$ 平凡的表示 $\pi$，正是需要证明 Ramanujan 猜想的表示，因此需要对所有 $\rho$ 证明 $L(s,\pi,\rho)$ 延拓到区域 $\operatorname{Re}s>1$。

无论如何，我们会得到一个类似于 (1.4) 的和，只是和中只出现假定为 Ramanujan 型的表示 $\pi$。这样我们就分离出了 $G$ 的 Ramanujan 型表示。接下来，对给定的 $\rho$，需要分离出那些 $L(s,\pi,\rho)$ 在 $s=1$ 或在临界线 $\operatorname{Re}s=1$ 上任意给定点有给定阶极点的表示。为分离出给定表示，需要利用 (1.4) 中 $f_v$（$v\in S$）的任意选择。

我们可以从迹公式出发，对大致任意的函数 $f=\prod_v f_v$，把 $\operatorname{tr}(\pi(f))$ 在所有自守表示 $\pi$ 上的和表达出来。起初（[L3]）似乎最好在 (1.4) 中放入的不是 (1.1) 中的 $L$ 函数（它们是乘积），而是它们的对数导数，因为对这些导数，留数是可加的。这导致类似
\[
\sum_{p\le X}\ln p
\]
的和，后者出现在素数定理的证明中。我们称它们为\textbf{算术和}。虽然我们稍后会转到几何和 (1.4)，我们先从算术和开始。

\srcpage{7}
若 $\pi$ 在 $S$ 外非分歧，则 (1.3) 的对数导数取负号等于
\begin{equation}\tag{1.5}
-\frac{(L_S)'(s,\pi,\rho)}{L_S(s,\pi,\rho)}
=-\sum_{v\notin S}\frac{L'_v(s,\pi,\rho)}{L_v(s,\pi,\rho)}
=\sum_{v\notin S}\ln N\p_v\sum_n\frac{\operatorname{tr}\rho^n(A(\pi_v))}{N\p_v^{ns}}.
\end{equation}
如果掌握了 Ramanujan 猜想，就可以把 (1.5) 写成和
\begin{equation}\tag{1.6}
\sum_v\ln N\p_v\frac{\operatorname{tr}\rho(A(\pi_v))}{N\p_v^s}
+\sum_{v\notin S}O\!\left(\frac{1}{N\p_v^{2}}\right).
\end{equation}
\yizhu{原文此处印作 $N\mathfrak p_v^{2}$，没有 $s$；按 (1.5) 的 Ramanujan 估计，$n\ge2$ 的项应为 $N\mathfrak p_v^{-2s}$，译文照原文排印，读者按 $2s$ 理解即可。}{依据：原书 p.7 页面图；对照 (1.5)}
然而这一结果在我们的起点上并不可用。它是一个目标。

如已解释，并不预期对单个表示 $\pi$ 计算 (1.5)。稳定迹公式给出的是
\begin{equation}\tag{1.7}
\sum_{v\notin S}\ln N\p_v\sum_n\left(\sum_{\pi^{\mathrm{st}}}m(\pi^{\mathrm{st}})
\prod_{v\in S}\operatorname{tr}\pi^{\mathrm{st}}_v(f_v)
\frac{\operatorname{tr}\rho^n(A(\pi_v))}{N\p_v^{ns}}\right),
\end{equation}
或
\begin{equation}\tag{1.8}
\sum_{v\notin S}\ln N\p_v\left(\sum_{\pi^{\mathrm{st}}}m(\pi^{\mathrm{st}})
\prod_{v\in S}\operatorname{tr}\pi^{\mathrm{st}}_v(f_v)
\frac{\operatorname{tr}\rho(A(\pi_v))}{N\p_v^{s}}\right),
\end{equation}
因为在位 $v$ 的 Hecke 代数中可以找到函数 $K^n_v$ 使得
\[
\operatorname{tr}\rho^n(A(\pi_v))=\operatorname{tr}\pi_v(K^n_v).
\]
\yizhu{原文 (1.7)、(1.8) 中，外层求和指标是 $v\notin S$，而 $\dfrac{\operatorname{tr}\rho^{n}(A(\pi_v))}{N\p_v^{ns}}$（(1.8) 中为 $\dfrac{\operatorname{tr}\rho(A(\pi_v))}{N\p_v^s}$）却印在 $\prod_{v\in S}$ 之内，与乘积指标 $v$ 冲突；按 (1.5) 与本节后文的用法，它应乘在乘积之外、对应外层位 $v$。译本照原文排版，特此说明。}{依据：原书 p.7；对照 (1.5)}

$G(F_v)$ 上的函数 $f_v$ 光滑紧支，除此之外任意。这使我们能在某些目的下分离出涉及的表示，更确切地说是 $L$-包，记作 $\pi^{\mathrm{st}}$，因为出现在稳定迹公式中的正是它们。重数 $m(\pi^{\mathrm{st}})$ 是\textbf{稳定重数}；这一概念的一般而精确的定义，无疑要随着稳定迹公式的建立才逐步给出。目前只在特殊情形下给出过；

\srcpage{8}
例如见 [LL]。无论如何，它将是一个可能为分数的数。我们提醒，$L_S(s,\pi,\rho)$ 只依赖于包含 $\pi$ 的 $L$-包。因此按定义它是稳定的。\yizhu{稳定重数与 $L$-包是 Arthur 内窥分类的基本对象：同一 $L$-包内的表示共享同一个稳定迹公式项，故稳定重数可以是分数。本库既有卡片把「稳定迹公式」与「$L$-包」列为 Arthur–Selberg 理论的入口。}{本库 lbook04.concept.ch8-arthur-selberg（p.188–191）；依据：原书 p.7–8}

迹公式把 (1.7) 表达为对（稳定）共轭类的和，困难在于分析它给出的等式。为我们的目的，最好先取 ${}^LG=\hat G\rtimes\operatorname{Gal}(K/F)$ 的一个不可约非平凡表示 $\rho$，其中 Galois 扩张足够大。例如，对 $G=\{1\}$，$\rho$ 是 Galois 群的任意有限维不可约复表示。目前看来，$L(s,\pi,\rho)$ 在直线 $\operatorname{Re}s=1$ 上的零点问题不应困扰我们。它们被认为不存在。我们已经指出，(1.7) 或 (1.8) 中 $\varphi$ 参数非平凡的表示 $\pi$ 的贡献可以递推处理，因为 $\varphi$ 的中心化子的维数小于 ${}^LG$ 的维数。最终会得到一个公式，其中出现的所有自守表示都是 Ramanujan 型，从而原则上使 (1.7)、(1.8) 中出现的 $L$ 函数在直线 $\operatorname{Re}s=1$ 右侧没有极点。\yizhu{原文此处语句有脱误：「les fonctions $L$ pour lesquelles les sommes (1.7) ou (1.8) n'ont pas de pôle…」，直译成中文是「使 (1.7) 或 (1.8) 中的和没有极点」，但「和」本身没有极点可言；按上下文应指其中出现的 $L$ 函数，译文据此译出。}{依据：原书 p.8；对照 (1.7)、(1.8)}把两个表达式之一乘以 $s-1$ 并令 $s\downarrow1$，就得到极点的阶。还有其他解析方法可从 (1.5) 或 (1.6) 中提取极点的阶，可能更有用；但这里重要的是，这些信息都蕴含在这两个表达式中。

若 $\pi$ 是 Ramanujan 型，人们预期——且有很好的理由——$L(s,\pi,\rho)$ 的对数导数在 $s=1$ 的极点的阶 $\mu(\pi)=\mu(\pi^{\mathrm{st}})$，等于平凡表示在 $\rho$ 限制到 ${}^{\lambda}H_\pi$ 中的重数；于是和
\begin{equation}\tag{1.9}
\sum_{\pi^{\mathrm{st}}}\mu(\pi)m(\pi^{\mathrm{st}})\prod_{v\in S}\operatorname{tr}\pi^{\mathrm{st}}_v(f_v)
\end{equation}
出现在迹公式中，原则上可以从中提取。若 $\rho$ 不可约且非平凡，则除少数满足 ${}^{\lambda}H_\pi\ne{}^LG$ 的 $\pi$ 外，这个重数为零。因此对大多数 $\pi$ 它为零。其余的 $\pi$ 来自群 ${}^{\lambda}H_\pi$，从而来自群 $H$；这些群的维数都小于 $G$ 的维数，因此原则上已被理解——但前提是函子性已经建立。利用它们的迹公式以及 $\rho$ 到 ${}^LH\to{}^LG$ 的限制 $\rho_H$，我们原则上可以把 (1.7) 计算为对 $H$ 的和。若平凡表示在 $\rho_H$ 中的重数为 $\mu(\rho_H)$，

\srcpage{9}
则和 (1.7) 应等于
\begin{equation}\tag{1.10}
\sum_H\mu(\rho_H)\sum_{\pi^{\mathrm{st}}_H}m(\pi^{\mathrm{st}}_H)
\prod_{v\in S}\operatorname{tr}\pi^{\mathrm{st}}_{H,v}(f_v).
\end{equation}
这个和蕴含在附属于群 ${}^{\lambda}H$ 的群 $H$ 的迹公式中，尽管必须考虑伴随包含链 ${}^LG\supset{}^{\lambda}H\supset\cdots$ 而来的麻烦。

刚才描述的东西会是对函子性的一个确认，但这不是我们所提议的。我们提议的是用同样的论证来验证函子性。这并不容易，因为迹公式中出现的对共轭类求和的分析并不容易（[L3]）。下一节的引理提示了另一种可能。虽然对数有重要的形式优点：重数 $\mu(\pi)$ 线性地出现，不必受无用因子的拖累，但从解析角度看，处理 $L$ 函数本身似乎更容易。除了分析本身的困难外，(1.9) 与 (1.10) 的等式中还隐藏着一个非常严重的困难，即对于嵌入 ${}^{\lambda}H\to{}^LG$ 的函子性。这里还涉及一个较小的困难，我们将在适当时候讨论。嵌入 ${}^LG$ 的群 ${}^{\lambda}H$ 本身并不是某个群 $H$ 的对偶群。它是一个群 ${}^LH$ 在某个可容许满射下的像。这无关紧要，但需要解释。关键的困难是，我们提出的策略要求我们原则上再次使用迹公式，来证明映射 ${}^LH\to{}^LG$ 的函子性。迄今为止我们不知道如何做到。我们一步步接近这个问题。下面解释策略。

定义 Hecke 算子 $K^{\rho,(n)}_v$，使得
\begin{equation}\tag{1.11}
\operatorname{tr}\bigl(\pi_v(K^{\rho,(n)}_v)\bigr)=\operatorname{tr}\rho^{(n)}(A(\pi_v)),
\end{equation}
其中 $\rho^{(n)}$ 是 $\rho$ 的 $n$ 次对称积。令
\[
L_v(s,\rho)=\sum_{n=0}^{\infty}q_v^{-ns}K^{\rho,(n)}_v.
\]
虽然我们随后引入的和与积在 $\operatorname{Re}s$ 充分大时收敛，但把它们当作 $t=q^{-s}$ 的形式幂级数处理有一些好处，其中 $N\p_v=q_v=q^{\deg v}$；或在数域情形，当作形式 Dirichlet 级数。

\srcpage{10}
显然
\begin{equation}\tag{1.12}
\operatorname{tr}\bigl(\pi_v(L_v(s,\rho))\bigr)=L_v(s,\pi_v,\rho)
\end{equation}
若 $\pi_v$ 非分歧，否则为 $0$。

我们在本文及后续文章中的目标，是开始讨论稳定迹公式的可能推论。我们准备略过一些基本问题，因为对我们起初处理的问题它们不造成困难。在系统化的理论中，称 $G(F_v)$ 的不可约表示 $\pi_v$ \textbf{非分歧}，若 $G_v$ 在 $F_v$ 上拟分裂、在某个非分歧扩张上分裂，且存在一个被给定超特殊子群固定的非零向量。称 $L$-包 $\pi^{\mathrm{st}}_v$ 非分歧，若其中每个元素对适当选择的这个超特殊子群都非分歧。从 $F$ 上的群 $G$——更确切地说，从定义它的一族方程——不仅得到每个 $F_v$ 上的群，而且对几乎所有有限位 $v$ 得到 $\mathcal{O}_v$ 上的群。子群 $G(\mathcal{O}_v)$ 几乎处处是超特殊的。我们选择 $S$ 时要求 $S$ 外如此。若选择另一族方程，则这样得到的子群 $G(\mathcal{O}_v)$ 几乎处处相同。作了这些选择后，等式 (1.11) 在 $S$ 外成立，但只对其中出现的 $L$-包 $\pi_v$ 的一个元素成立。这个元素与 Hecke 代数由 $G(\mathcal{O}_v)$ 的选择固定。对包中其他元素，$\operatorname{tr}(\pi_v(K^{(n)}_v))=0$。因此 (1.11) 与 (1.12) 的精确形式是
\[
\operatorname{tr}\bigl(\pi^{\mathrm{st}}_v(K^{(n)}_v)\bigr)=\operatorname{tr}\rho^{(n)}(A(\pi_v))
\quad\text{与}\quad
\operatorname{tr}\bigl(\pi^{\mathrm{st}}_v(L_v(s,\rho))\bigr)=L_v(s,\pi_v,\rho).
\]
分解式
\[
L_S(s,\pi,\rho)\prod_{v\in S}\operatorname{tr}(\pi^{\mathrm{st}}_v(f_v))
=\prod_{v\notin S}\operatorname{tr}\bigl(\pi_v(L_v(s,\rho))\bigr)
\prod_{v\in S}\operatorname{tr}(\pi^{\mathrm{st}}_v(f_v))
\]
蕴含和
\begin{equation}\tag{1.13}
\sum_{\pi^{\mathrm{st}}}m(\pi^{\mathrm{st}})L_S(s,\pi,\rho)\prod_{v\in S}\operatorname{tr}(\pi^{\mathrm{st}}_v(f_v))
\end{equation}
由稳定迹公式对函数
\begin{equation}\tag{1.14}
f=\bigotimes_{v\notin S}L_v(s,\rho)\bigotimes_{v\in S}f_v
\end{equation}
给出。这些公式中混合了加法与乘法。然而可以希望仍以同样的方式

\srcpage{11}
递推地排除非 Ramanujan 型表示的贡献。剩下的和将等于：对初始群 $G$ 的迹公式给出的和，与对 ${}^LG$ 中 $\operatorname{SL}(2)$ 的非平凡嵌入所联系的群的和，二者之差。在第五节我们处理最简单的情形，即 $\operatorname{SL}(2)$ 到 ${}^LG$ 的主嵌入，并说明如何在 $G$ 的稳定迹公式中辨认附属于它的表示。本文中我们不会走得更远；与其说我们对结果有信心，不如说我们对新原理的优雅与简单有信心。我们希望在不远的将来回到这些问题，即从迹公式中提取其他非 Ramanujan 型表示——从而联系于既非平凡也非主嵌入的 $\varphi$——的贡献。

完成这些之后，从稳定迹公式将得到某个表达式，记作 $\Xi(s)$，它类似于 (1.13) 的和，但原则上只涉及 Ramanujan 型表示。对这些表示，$L_S(s,\pi,\rho)$ 在 $\operatorname{Re}s>1$ 中应既无极点也无零点，人们希望能从该表达式得出这一结论。无论如何，对这个表达式有贡献的 $\pi^{\mathrm{st}}$ 将分属不同类型：来自 ${}^{\lambda}H$ 且 $\rho$ 限制到 ${}^{\lambda}H$ 不含平凡表示的，以及其余的。对给定的 $\rho$，前者不影响 $\Xi(s)$ 在 $s=1$ 的极点，可以弃之不顾。

通过转移来自群 $H$、且 $\rho$ 与同态 ${}^{\lambda}H\to{}^LG$ 的复合含平凡表示的那些类，是唯一对主部有贡献的类。这些群 $H$ 的维数小于 $G$，因此原则上已被理解。于是可以从它们的稳定迹公式计算它们对 $\Xi(s)$ 在 $s=1$ 的主部的贡献。把所有 $H$ 的这些主部相加，并与 $\Xi(s)$ 的主部比较，应能得到一个再次确认函子性的等式，甚至更好，导致一般情形的证明。我们在第五节以谦逊的方式着手这些计算，但真正的工作仍待完成。

这一策略的核心有一个尚未完全解决的困难，但我们认为并不严重。我们使用的是 $L$ 函数而非其对数导数。因此，若 $\pi_H$ 不是来自维数小于 $H$ 的群 $H'$ 的函子性像，且 $\rho$ 限制到 ${}^{\lambda}H$ 是

\srcpage{12}
$m$ 倍平凡表示与一个不含平凡表示的表示 $\tau$ 的和，则原则上
\[
L_S(s,\pi,\rho)=\zeta_F^m(s)L_S(s,\pi,\tau),
\]
其中第二个因子在 $s=1$ 既无零点也无极点。然而它对 $s=1$ 的主部有影响。相反，稳定迹公式是对 $\pi$ 的稳定类的和（！），并不把附属于不同 $\pi^{\mathrm{st}}$ 类的 $L$ 函数乘在一起。

主要的困难将是分析 (1.13) 的和在去掉所有 $\varphi$ 非平凡项之后的渐近行为。我们尚未开始，但第 3、4、5 节是为这一目的作准备。
